\begin{pastexam}[label={px:ust-114-ee-ds-14}]{臺聯大電機類 114 學年度 第 14 題}
（單選；2 分，答錯不倒扣）

Answer questions 14-16 according to the following C code.
\begin{minted}[fontsize=\small]{c}
struct Node {
    int data;
    struct Node* next;
};
int funcA(struct Node* head) {
    struct Node *slow = head, *fast = head;
    while (slow && fast && fast->next) {
        slow = slow->next;
        fast = fast->next->next;
        if (slow == fast) {
            return 1;
        }
    }
    return 0;
}
\end{minted}

What is the purpose of funcA()?

(A)~Find repeated numbers in a linked list \qquad (B)~Find the unique numbers in a linked list\\
(C)~Find a loop in a linked list \qquad (D)~Reverse a linked list
\end{pastexam}

\begin{pastsoln}
答案是 (C)。

這是 Floyd 的快慢指標（龜兔賽跑）判環法（資結書第~5.6~節）。 \texttt{slow} 每次走一步，
\texttt{fast} 每次走兩步。
\begin{itemize}
\item 沒有環：\texttt{fast} 會先走到串列尾端，\texttt{fast} 或 \texttt{fast->next} 變成
\texttt{NULL}，迴圈結束，回傳 $0$。 兩個指標只往前走，\texttt{slow} 落後 \texttt{fast}，不會相等。
\item 有環：兩個指標最後都進入環裡。 之後每走一輪，\texttt{fast} 和 \texttt{slow} 沿環方向的距離
少 $1$，距離最多是環長減一，所以在環長以內的輪數裡一定相遇，回傳 $1$。
\end{itemize}
所以 \texttt{funcA} 判斷串列有沒有環（loop）。
\begin{itemize}
\item (A)(B)：程式完全沒有讀 \texttt{data}，只比較指標，和數值重不重複無關。 錯。
\item (C)：正確。
\item (D)：沒有改任何 \texttt{next} 指標，不會反轉串列。 錯。
\end{itemize}
\end{pastsoln}
